\chapter{Technology Selection and Platform Constraints}
\label{chap:technologies-and-platform-constraints}

This chapter evaluates the technology, framework, and deployment decisions that define the platform for the prototype, rather than presenting the complete final system design. It is organised by system layer: rendering and 3D graphics, audio capture and processing, machine learning deployment, and frontend application architecture. For each layer, the selected tools are presented, alternatives are compared, and the rationale for selection is justified. The chapter concludes by deriving the platform constraints these decisions impose, which establish the design space developed methodologically and architecturally in Chapter 4.

\section{The Web as a Platform for Real-Time Audio-Visual Applications}

The decision to build this system as a web application is motivated by three considerations: accessibility, portability, and the maturity of the modern browser as an application runtime. A web application requires no installation, runs in any modern browser regardless of operating system, and can be shared with collaborators or evaluators by distributing a URL rather than a platform-specific binary. For a research prototype intended to demonstrate feasibility, this low barrier to entry is a meaningful advantage.

These benefits depend on the browser providing the capabilities required for real-time audio-visual processing. The modern browser exposes the necessary APIs: WebGL for hardware-accelerated 3D rendering, the Web Audio API for low-latency audio capture and spectral analysis, WebAssembly for near-native computational performance, and Web Workers for background-thread execution. Collectively, these APIs make interactive audio-visual applications viable within the browser sandbox.

Alternative approaches were considered. A native C++/OpenGL application would offer unrestricted GPU access and lower audio latency, but would require platform-specific builds and significantly more complex distribution. Electron inherits the browser's resource overhead while losing the zero-install advantage. Game engines such as Unity or Unreal Engine provide powerful rendering and audio capabilities, but introduce a heavy runtime and licensing considerations disproportionate for a stage-lighting research prototype. The web platform, complemented by a local server for ML inference (Section 3.4), provides the best balance of accessibility, development speed, and technical capability for this project.

\section{3D Rendering: WebGL and Three.js}

The visual core of this application is a real-time 3D rendering of a virtual stage with lighting fixtures. WebGL provides the browser's interface to the GPU: a JavaScript API based on OpenGL ES 2.0 that exposes a programmable rendering pipeline without plugins \cite{threejs-docs}. However, WebGL operates at a low level of abstraction, requiring manual management of vertex buffers, GLSL shader programs, and the rendering state machine, which makes it impractical for application-level development.

Three.js addresses this by providing a high-level, scene-graph-based abstraction over WebGL. As the dominant web 3D library \cite{threejs-github,npm-three}, it provides the building blocks this project requires: scene and camera management, a renderer, and a hierarchy of meshes, materials, and lights. Most relevant to this project is the SpotLight primitive, which exposes parameters for colour, intensity, cone angle, penumbra, distance, decay, and target direction. These map directly to the physical characteristics of stage lighting fixtures (PAR cans, moving heads), making SpotLight a natural abstraction for the virtual stage \cite{threejs-docs}. Three.js also supports shadow mapping, post-processing pipelines (bloom, colour grading), and fog simulation, providing the visual vocabulary needed for a convincing stage environment.

The principal alternative, Babylon.js, offers a more comprehensive feature set (built-in physics, visual editor, native WebXR) but produces significantly larger bundles, often exceeding 1 MB even with tree-shaking \cite{babylonjs-forum-bundle-size}, compared to Three.js's approximately 155 KB gzipped \cite{bundlephobia-three}. Since this project requires only lighting, camera controls, and scene composition, Three.js's lighter footprint, larger community, and mature React integration (Section 3.5) make it the better fit.

WebGPU, a next-generation API superseding WebGL, has reached production readiness across all major browsers as of late 2025 \cite{webdev-webgpu-major-browsers-2025,gpuweb-implementation-status}. Three.js already supports a WebGPU renderer alongside its WebGL renderer, so choosing Three.js preserves a clear upgrade path without requiring architectural changes. This project targets the WebGL renderer for maximum compatibility.

\begin{figure}[htbp]
    \centering
    \includegraphics[width=\textwidth]{figures/chapters/3/threejs_scene_graph_and_lighting_model.pdf}
    \caption{Three.js scene graph and lighting model used in this project. The diagram shows a \texttt{Scene} hierarchy with \texttt{PerspectiveCamera}, lighting nodes (\texttt{SpotLight}, \texttt{PointLight}, and \texttt{SpotTarget}), mesh groups for stage elements, and the \texttt{WebGLRenderer} consuming both scene and camera inputs to produce the rendered frame. SpotLight parameters correspond to stage-lighting controls: colour (hue), intensity (value), angle (beam width), penumbra (edge softness), and target (direction).}
\end{figure}

\section{Audio Capture and Processing: The Web Audio API}

The Web Audio API is a W3C specification that provides a graph-based audio processing architecture in the browser \cite{w3c-webaudio-1-1}. Unlike HTML \texttt{} elements, it exposes a modular system of audio nodes connected into a directed processing graph, enabling real-time analysis, transformation, and routing with low latency \cite{mdn-web-audio-api}. An AudioContext manages the graph and processes audio in render quanta of 128 samples (approximately 2.67 ms at 48 kHz) \cite{mdn-audiocontext}. This application supports two audio input modes that produce an identical real-time stream from the perspective of downstream processing. For microphone input, \texttt{getUserMedia()} provides a MediaStream wrapped as a MediaStreamAudioSourceNode \cite{mdn-getusermedia}. For file-based input, the uploaded audio is decoded into an AudioBuffer and played back through an AudioBufferSourceNode. Crucially, both modes feed the same processing graph frame by frame as audio plays, rather than pre-analysing the entire file. This unified streaming design ensures the system exercises identical real-time constraints regardless of source.

Two processing nodes are central to the pipeline. The AnalyserNode performs a real-time FFT on each audio block, exposing frequency-domain and time-domain data without modifying the audio stream \cite{mdn-analysernode-fftsize}. It provides spectral data for the frontend visualisation and can supply initial frequency features to complement server-side analysis. The AudioWorklet runs a user-defined processor in a dedicated audio thread, separate from the main JavaScript thread, ensuring audio processing is not interrupted by UI activity \cite{mdn-audioworklet-guide}. In this project, the AudioWorklet extracts raw PCM audio chunks and transmits them to the server via WebSocket for consumption by the beat tracker and Skip-BART model.

The API imposes several design-relevant constraints. Browsers require a user gesture before an AudioContext can be activated (autoplay policy). The total audio output latency, the sum of \texttt{baseLatency} and \texttt{outputLatency}, typically ranges from 10 to 50 milliseconds on standard hardware \cite{mdn-audiocontext}. Communication between the AudioWorklet thread and the main thread requires message passing via \texttt{port.postMessage()} or SharedArrayBuffers, since the AudioWorklet runs in an isolated scope.

\begin{figure}[htbp]
    \centering
    \includegraphics[width=0.85\textwidth]{figures/chapters/3/web_audio_api_routing_graph.pdf}
    \caption{Representative Web Audio API routing configuration enabled by the selected browser stack. Microphone and audio-file inputs converge into a shared chain (\texttt{GainNode} → \texttt{AnalyserNode} → \texttt{AudioWorkletNode} → \texttt{AudioDestinationNode}), illustrating support for simultaneous playback, spectral analysis, and low-latency audio streaming to the inference backend.}
\end{figure}

\section{Machine Learning Deployment}
\label{sec:ch3-machine-learning-deployment}

The system depends on two categories of ML models from Chapter 2: a beat tracker (BeatNet+, BEAST, or PLP) for real-time rhythm analysis, and Skip-BART for generative lighting control. Both were developed in Python using PyTorch. Deploying them alongside a web application requires bridging the Python ML ecosystem and the browser runtime.

\subsection{Client-Side Inference Runtimes}

Two runtimes support ML inference directly in the browser. ONNX Runtime Web runs models in ONNX format using WebAssembly (CPU) or WebGPU (GPU) backends, with PyTorch models exportable via \texttt{torch.onnx.export()} \cite{onnxruntime-web-tutorial}. TensorFlow.js provides an alternative with WebGL, WebAssembly, and WebGPU backends, though models originating in PyTorch require a multi-step conversion path (PyTorch → ONNX → TensorFlow.js) \cite{tensorflowjs-platform-environment}. However, client-side inference faces several constraints critical to this project. Skip-BART's 240 million parameters would require roughly 960 MB in FP32 (480 MB in FP16), approaching the 1 to 4 GB per-tab browser memory limits. Model conversion can fail with custom operators or dynamic control flow. The beat tracking models rely on Python-specific post-processing (Dynamic Bayesian Networks, particle filters, RSTC sampling) that would require reimplementation in JavaScript. Finally, browser GPU inference via WebGL/WebGPU remains less efficient than native CUDA.

\subsection{Server-Side Inference}

The alternative is running models natively in Python/PyTorch on a local server, which eliminates model conversion entirely. The original research implementations can execute with their published weights and post-processing logic, and a CUDA-capable GPU can be leveraged directly.

FastAPI is selected as the server framework because it provides native WebSocket support and an asynchronous architecture suitable for real-time streaming \cite{fastapi-docs}. This choice allows the Python model stack to remain unchanged while supporting a low-latency browser-to-server streaming loop.

\subsection{Real-Time Communication}
\label{sec:ch3-realtime-communication}

The hybrid architecture requires a persistent, bidirectional, low-overhead communication channel. The WebSocket protocol satisfies all three requirements: after an initial HTTP handshake, it maintains a full-duplex connection with minimal per-message framing overhead (2 to 14 bytes versus HTTP's 200 to 800 bytes of headers per request) \cite{postman-websockets-vs-http}. On localhost, round-trip latency is sub-millisecond, making the network hop negligible relative to model inference time.

\subsection{Justification for the Hybrid Architecture}

Figure 3.3 summarises the deployment alternatives evaluated for this thesis. The decision to adopt server-side inference connected to the browser via WebSockets is driven by four considerations: (1) Skip-BART's 240 million parameters approach or exceed browser memory limits; (2) the Python-specific post-processing logic has no direct ONNX/TensorFlow.js equivalent, and reimplementing it in JavaScript would be a significant effort tangential to the thesis contribution; (3) native PyTorch with CUDA acceleration provides substantially faster inference, particularly for Skip-BART's autoregressive decoding; and (4) the clean separation into a rendering frontend and an analysis backend simplifies development and mirrors standard production ML architectures. The trade-off is an operational dependency on a running server, which is acceptable for a thesis demonstration context and is carried forward as a design constraint in Chapter 4.

\begin{figure}[htbp]
    \centering
    \includegraphics[width=\textwidth]{figures/chapters/3/client_vs_server_inference_architecture.pdf}
    \caption{Client-side versus server-side inference architecture comparison. The left panel shows a pure browser pipeline (\texttt{Web Audio API}, in-browser model inference, post-processing, and Three.js rendering) and highlights its constraints: memory limits, lack of CUDA acceleration, model-conversion overhead, and post-processing reimplementation cost. The right panel shows the chosen hybrid architecture, where audio capture and rendering remain in the browser while native PyTorch beat tracking and Skip-BART inference run on a Python server, connected through bidirectional WebSockets.}
\end{figure}

\section{Frontend Framework Selection: React and React Three Fiber}

The frontend must support three concurrent concerns: rendering an interactive 3D stage, presenting control UI, and integrating high-frequency model outputs. The application is written in TypeScript for compile-time type safety across data contracts (audio features, WebSocket messages, and scene parameters) \cite{typescript-official}. React provides the declarative component model for the UI layer (audio source selection, playback controls, fixture configuration). The 3D scene is implemented with React Three Fiber (R3F), a React renderer for Three.js that expresses Three.js objects as JSX components without abstracting away or adding overhead to the underlying library \cite{r3f-github}. R3F's companion library, Drei \cite{drei-github} provides ready-made components directly applicable to this project: OrbitControls for camera navigation, Environment for ambient lighting, and a SpotLight helper with volumetric cone visualisation.

State management uses Zustand, a subscription-based store from the same Poimandres collective \cite{zustand-github}. Compared with heavier alternatives such as Redux-based setups, Zustand offers lower boilerplate and straightforward selective subscriptions, which is advantageous for real-time interfaces that mix UI concerns with rapidly changing visual parameters.

Alternative frameworks were evaluated. Vue with TresJS and Svelte with Threlte both offer declarative Three.js integration, with Svelte's compile-time approach eliminating framework runtime overhead. However, React Three Fiber's substantially larger ecosystem (companion libraries, community documentation, solved integration patterns) was the deciding factor for a thesis project where development speed is critical.

The selected React + R3F + Zustand stack therefore provides the required capability envelope for responsive UI control and high-frequency 3D updates in a browser setting. The detailed module responsibilities and runtime data-flow specification are presented in Chapter 4, where the full system design is developed.

\section{Platform Constraints and Design Implications}

The technologies presented in this chapter impose constraints that delimit the system design space.

JavaScript's single-threaded execution model means the main thread is shared between React rendering, DOM updates, and event handling. The AudioWorklet and Web Workers provide partial mitigation, but coordination still requires asynchronous message passing and explicit separation of UI work from time-sensitive audio/visual processing.

Browser memory limits (1 to 4 GB per tab) were a primary factor in adopting server-side ML inference. Skip-BART's 240 million parameters alongside the Three.js scene, React state, and audio buffers would leave insufficient headroom. WebGL's higher abstraction level compared to native APIs (no compute shaders, less efficient GPU memory management) is unlikely to be a bottleneck for this project's modest stage scene, but constrains future extensions. WebGPU addresses this prospectively (Section 3.2).

The Web Audio API's latency floor (10 to 50 ms total from \texttt{baseLatency} plus \texttt{outputLatency}) is additive with WebSocket round-trip time and server inference time. The perceptual thresholds established in Chapter 2 indicate that total pipeline latencies below approximately 100 ms are perceived as synchronous, requiring the latency budget to be managed across the entire chain. Cross-browser differences in WebGL extensions, AudioWorklet behaviour, and WebSocket handling led to targeting Chrome/Chromium as the primary runtime.

Finally, the hybrid architecture requires a running Python server, preventing standalone web deployment. For the thesis demonstration this is acceptable; containerisation (Docker) or cloud hosting could extend accessibility in future work.

These constraints define the system's design space. On that basis, Chapter 4 presents the methodological approach, derives the key requirements, and specifies the resulting system design in detail, including module boundaries, runtime data flow, and major design decisions.
